\documentclass[10pt, a5paper]{article}
\usepackage{pdfpages}
\usepackage{parallel}
\usepackage[T2A]{fontenc}
\usepackage{ucs}
\usepackage[utf8x]{inputenc}
\usepackage[polish,english,russian]{babel}
\usepackage{hyperref}
\usepackage{rotating}
\usepackage[inner=2cm,top=1.8cm,outer=2cm,bottom=2.3cm,nohead]{geometry}
\usepackage{listings}
\usepackage{graphicx}
\usepackage{wrapfig}
\usepackage{longtable}
\usepackage{indentfirst}
\usepackage{array}
\newcolumntype{P}[1]{>{\raggedright\arraybackslash}p{#1}}
\frenchspacing
\usepackage{fixltx2e} %text sub- and superscripts
\usepackage{icomma} % коскі ў матэматычным рэжыме
\PreloadUnicodePage{4}

\newcommand{\longpage}{\enlargethispage{\baselineskip}}
\newcommand{\shortpage}{\enlargethispage{-\baselineskip}}

\def\switchlang#1{\expandafter\csname switchlang#1\endcsname}
\def\switchlangbe{
\let\saverefname=\refname%
\def\refname{Літаратура}%
\def\figurename{Іл.}%
}
\def\switchlangen{
\let\saverefname=\refname%
\def\refname{References}%
\def\figurename{Fig.}%
}
\def\switchlangru{
\let\saverefname=\refname%
\let\savefigurename=\figurename%
\def\refname{Литература}%
\def\figurename{Рис.}%
}

\hyphenation{admi-ni-stra-tive}
\hyphenation{ex-pe-ri-ence}
\hyphenation{fle-xi-bi-li-ty}
\hyphenation{Py-thon}
\hyphenation{ma-the-ma-ti-cal}
\hyphenation{re-ported}
\hyphenation{imp-le-menta-tions}
\hyphenation{pro-vides}
\hyphenation{en-gi-neering}
\hyphenation{com-pa-ti-bi-li-ty}
\hyphenation{im-pos-sible}
\hyphenation{desk-top}
\hyphenation{elec-tro-nic}
\hyphenation{com-pa-ny}
\hyphenation{de-ve-lop-ment}
\hyphenation{de-ve-loping}
\hyphenation{de-ve-lop}
\hyphenation{da-ta-ba-se}
\hyphenation{plat-forms}
\hyphenation{or-ga-ni-za-tion}
\hyphenation{pro-gramming}
\hyphenation{in-stru-ments}
\hyphenation{Li-nux}
\hyphenation{sour-ce}
\hyphenation{en-vi-ron-ment}
\hyphenation{Te-le-pathy}
\hyphenation{Li-nux-ov-ka}
\hyphenation{Open-BSD}
\hyphenation{Free-BSD}
\hyphenation{men-ti-on-ed}
\hyphenation{app-li-ca-tion}

\def\progref!#1!{\texttt{#1}}
\renewcommand{\arraystretch}{2} %Іначай формулы ў матрыцы зліпаюцца з лініямі
\usepackage{array}

\def\interview #1 (#2), #3, #4, #5\par{

\section[#1, #3, #4]{#1 -- #3, #4}
\def\qname{LVEE}
\def\aname{#1}
\def\q ##1\par{{\noindent \bf \qname: ##1 }\par}
\def\a{{\noindent \bf \aname: } \def\qname{L}\def\aname{#2}}
}

\def\interview* #1 (#2), #3, #4, #5\par{

\section*{#1\\{\small\rm #3, #4. #5}}

\def\qname{LVEE}
\def\aname{#1}
\def\q ##1\par{{\noindent \bf \qname: ##1 }\par}
\def\a{{\noindent \bf \aname: } \def\qname{L}\def\aname{#2}}
}

\begin{document}
\title{Голос спонсора: Conjur. Authorization in the cloud}
%\author{}
\date{}
\maketitle

\begin{figure}[h!]
\centering
\includegraphics{74_conjur-logo-1.pdf}
\end{figure}

Conjur is a cloud-native platform for directory services, authoriza\-tion, and audit for development and operations teams and their entire infrastructure. 

With 100\% API coverage and a scalable, easily deployable, high-availability architecture, Conjur reduces the time, cost, and complexity associated with building authorization management via complex home\-grown scripts and configuration management tools. Sometimes referred to as “Active Directory for the cloud”, Conjur runs in either a virtual machine or container, and works alongside a wide range identity and access management (IAM) solutions to solve access and authorization challenges: machine-to-machine permissions, deployment and access \linebreak rights to 
sensitive systems, and the auditing required to meet complian\-ce requirements. 

“Conjur is more than just an outstanding platform for managing authorization as a service: they are a partner that we can innovate with, and an essential part of achieving our vision and solving the hard problems that arise as we move our infrastructure into the cloud technologies that will drive our business forward.” 
- Mike Kail, VP of Operations, Netflix 

\section*{Why Conjur? }
Built with system administrators, DevOps professionals, and cloud architects in mind, Conjur natively supports a wide range of potential use cases. Some of the most common reasons why organizations select Conjur over building in-house solutions include: 
\begin{itemize}
\item {\bf Compliance}: auditable enforcement of both organizational and regulatory policies and rules, via existing reporting systems (e.g., SIEM platforms, SumoLogic, Splunk) 
\item {\bf Risk Management}: reduction of the attack surface for sensitive data (credentials, SSL/SSH keys and certificates, secret s, etc.)
\item {\bf DevOps Optimization}: integration of security \& controls in upstream development and operations work, ensuring consistency across all production systems. 
\item {\bf Access Intelligence}: unified control of identity (human \& machi\-ne) and permissions across entire infrastructure (bare metal, pri\-vate, public, cloud) helps prevent failures, without relying on legacy systems for policy governance and enforcement 
\end{itemize}

\section*{What is Conjur?}


Conjur is designed with two primary goals: (1) to provide an easy-to-use, well documented, extensible authorization system, and (2) to allow for non-disruptive integration into organizati\-onal workflows. 

This approach has resulted in a series of implementation benefits: 
\begin{itemize}
\item {\bf Ease of Integration}: Conjur is built with a DevOps optimized UX (CLI-driven), with 100\% API access from a variety of langua\-ges, and can be deployed as a Linux virtual appliance, or into bare-metal systems 
\item {\bf No Vendor Lock-In}: Conjur has been used in conjunction with all leading configuration management and DevOps tools (Puppet, Chef, Salt, Docker, etc.) 
\item {\bf HA Design}: Conjur’s system has been designed with high availa\-bility and fault tolerance, with a distributed configuration archi\-tecture and full fail-over/redundancy capabilities 
\item {\bf Minimal Dependencies}: Conjur is self-contained, requiring no external database or server  
\item {\bf Secure}: All communication is fully encrypted, following industry best practices for managing data in use and at rest; additionally, Conjur has invested in a rigorous third-party security analysis (results expected to be published Q4 FY14)
\end{itemize}

Learn more at \url{http://conjur.net}

\end{document}


